\abstract{РЕФЕРАТ}

Объем работы равен \formbytotal{lastpage}{страниц}{е}{ам}{ам}. Работа содержит \formbytotal{figurecnt}{иллюстраци}{ю}{и}{й}, \formbytotal{tablecnt}{таблиц}{у}{ы}{}, \arabic{bibcount} библиографических источников и \formbytotal{числоПлакатов}{лист}{}{а}{ов} графического материала. Количество приложений -- 2. Графический материал представлен в приложении А. Фрагменты исходного кода представлены в приложении Б.

Перечень ключевых слов: компилятор, транслятор, лексер, парсер, код, ассемблер, виртуальная машина, байт-код, инструкция, регистр, кадр активации, стек, функциональное программирование, Common Lisp, S-выражение, макросы, рекурсия, оптимизация.

Объектом разработки является программная система компилятора языка программирования Common Lisp.

Целью выпускной квалификационной работы является создание программной системы компилятора языка программирования Common Lisp.

В ходе работы были выделены основные этапы компиляции и реализованы соответствующие модули: преобразование исходного S-выражения в промежуточную форму с упрощёнными операциями, оптимизация этой формы, генерация кода ассемблера и трансляция в байт-код виртуальной машины. Оптимизатор проходит промежуточное дерево в глубину.

Виртуальная машина имеет следующие регистры: программный счётчик для последовательного выполнения инструкций и возможности перехода, аккумулятор для промежуточных значений, указатели вершины стека общего назначения, текущего кадра окружения и вершины стека нелокальных переходов. Константы и глобальные переменные хранятся в массивах и загружаются в аккумулятор по индексу. Перед вызовом функции создаётся новый кадр активации, наследующий окружение объявления функция, и в этот кадр активации заносятся значения аргументов через стек.

При разработке компилятора использовался язык программирования Common Lisp.

Разработанный компилятор был успешно внедрен в os-pi.

\selectlanguage{english}
\abstract{ABSTRACT}
  
The volume of work is \formbytotal{lastpage}{page}{}{s}{s}. The work contains \formbytotal{figurecnt}{illustration}{}{s}{s}, \formbytotal{tablecnt}{table}{}{s}{s}, \arabic{bibcount} bibliographic sources and \formbytotal{числоПлакатов}{sheet}{}{s}{s} of graphic material. The number of applications is 2. The graphic material is presented in annex A. The source code fragments are presented in annex B.

List of keywords: compiler, translator, lexer, parser, code, assembler, virtual machine, bytecode, instruction, register, activation frame, stack, functional programming, Common Lisp, S-expression, macros, recursion, optimization.

The object of the development is a software system for a compiler of the Common Lisp programming language.

The purpose of the final qualifying work is the creation of a software system for a compiler of the Common Lisp programming language.

In the course of the work, the main stages of compilation were identified and the corresponding modules were implemented: conversion of the source S-expression into an intermediate form with simplified operations, optimization of this form, generation of assembler code, and translation into virtual machine bytecode. The optimizer traverses the intermediate tree depth-first.

The virtual machine has the following registers: a program counter for sequential instruction execution and branching capability, an accumulator for intermediate values, a general-purpose stack top pointer, a current environment frame pointer, and a non-local jump stack top pointer. Constants and global variables are stored in arrays and loaded into the accumulator by index. Before a function call, a new activation frame is created, inheriting the function's declaration environment, and argument values are placed into this activation frame via the stack.

The Common Lisp programming language was used in the development of the compiler.

The developed compiler was successfully deployed in os-pi.
\selectlanguage{russian}
